
\chapter*{Chercher avant de consulter}
\addcontentsline{toc}{chapter}{Chercher avant de consulter}
Ce recueil contient \textbf{180 exercices}, six par chapitre, et \textbf{huit
problèmes de synthèse}. Les numéros correspondent exactement à ceux des corrigés.
Commencer par les calculs et situations directes, puis travailler les preuves
et les contre-exemples. Les exercices des chapitres 15 et 25 à 30 sont des
compléments ; certains exercices des autres chapitres prolongent aussi le socle.

Pour une réponse rédigée, préciser le domaine, les hypothèses et l'unité,
expliquer le choix de la méthode, puis vérifier le résultat. Une valeur approchée
ne remplace une valeur exacte que si la question le demande.
Un résultat informatique doit être accompagné d'un raisonnement sur ce qu'il
vérifie et ce qu'il ne prouve pas.

Les problèmes de synthèse croisent plusieurs chapitres. Ils peuvent servir
à une recherche personnelle, à une présentation orale ou à un travail de groupe.
Pour progresser vers le supérieur, reprendre quelques jours plus tard une preuve
sans son texte est plus formateur que relire immédiatement une solution.
